\cotefile{157-5}
\scanpage{20}
\mspage{7}

\diagram{\begin{tikzcd}[ampersand replacement=\&,column sep=large]
 X'_{i_0}\arrow[r]\arrow[d,no head,"\alpha_{i_0}"']\arrow[dr,phantom,"\text{cocart}" description]
 \& X_{i_0}\arrow[r]\arrow[d,"\beta_{i_0}"']\arrow[dr,phantom,"\text{\unsure{cocart}}" description]
 \& X_{\unsure{\infty}}\arrow[d]\\
 Y'_{i_0}\arrow[r,no head] \& \widetilde Y_{i_0}\arrow[r]\arrow[d,"\widetilde u_{i_0}"']\arrow[dr,phantom,"\text{cocart}" description]
 \& Z'_{i_0}\arrow[d,"\unsure{v|u}_{i_0}"]\\
 \& Y_{i_0}\arrow[r,no head] \& Z_{i_0}\,.
\end{tikzcd}}

On a $\widetilde u_{i_0}\in C$ par hypothèse, donc $v_{i_0}\in C$, ce qui achève la démonstration.

\uline{Proposition 2} $(\unsure{\mathcal M|\mathscr M},C)$ et le cas échéant $W$, comme dans Prop.~1 et son corollaire, $I$ ens ordonné localement fini, i.e tel que $\forall i_0\in I$, $I_{\unsure{/|\leq}i_0}$ soit fini. Soit\deleted{\unsure{ent|s}} \deleted{\unsure{y|z}}

\[
 u_\bullet:X_\bullet\longrightarrow Y_\bullet
\]

flèche dans $ {}^{*}\mathcal M(I^\circ)_C={}$\inserted{\unsure{ss cat.} \deleted{\unsure{p}}pleine des} systèmes inductifs sur $I$, \unsure{à valeurs|internes} dans $\mathcal M$, tels
% slanted
\margin{en fait, l’hypothèse que $X_\bullet,Y_\bullet$ dans $\mathcal M(I^\circ)_C$ est inutile, \unsure{tant} dans la prop.~2 que dans les corollaires} que $\varnothing_{\mathcal M(I^\circ)}$\inserted{objet initial de $\mathcal M(I^\circ)$} $\mathbin{\unsure{-|\to}} Z_\bullet$ soit dans $C_I$ (\unsure{défini|définie} dans prop.~1), i.e. tel que $\forall i_0\in I$, la flèche $Z'_{i_0}=\underset{i<i_0}{\underrightarrow{\lim}}\,Z_i\longrightarrow Z_{i_0}$ soit dans $C$, et supposons $u_\bullet\in C_I$. Alors

a) $\underrightarrow{\lim}_{\unsure{I}}u_\bullet:\deleted{\unsure{X}}X_I\overset{\text{déf}}{=}\underset{i\in I}{\underrightarrow{\lim}}\,X_i\longrightarrow Y_I\overset{\text{déf}}{=}\underset{i\in I}{\underrightarrow{\lim}}\,Y_i$ est dans $C$,
% archive sheet number: 19

\scanpage{21}
\mspage{8}

et b) $\forall i_0\in I$, \inserted{\unsure{nos}} $\alpha_{i_0},\beta_{i_0}\in C$ dans le diagramme \unsure{de|ci} p.~1.

\deleted{\uline{\unsure{Corollaire}}} De plus (supposant donné $W$ comme dans prop.~1 cor.~1), \inserted{les conditions suivantes} \deleted{on a $u_\bullet\in W_{I^\circ}$ (i.e. $u_{i_0}\in W$ $\forall i_0\in I$) alors $\unsure{u_I}\in TC=C\cap W$ et pour \deleted{\unsure{tout}} $i_0\in I$, on a $\widetilde u_{i_0}\in TC$ \unsure{donc}} sur $u_\bullet$ entre $X_\bullet,Y_\bullet$ $C$-cofibrants sont équivalentes

\[
\begin{aligned}
\deleted{\unsure{2}}1^\circ)&\quad u_\bullet\in(TC)_I,\quad\text{i.e. les }\widetilde u_{i_0}\in TC\quad(\forall i_0\in I)\\
\deleted{\unsure{3}}2^\circ)&\quad u_\bullet\in T_I\cap W_{I^\circ},\quad\text{i.e. }\forall i_0\in I,\quad\widetilde u_{i_0}\in C,\ u_{i_0}\in W,
\end{aligned}
\]

\deleted{\unsure{Alors}} $3^\circ)$ Pour tout $i_0\in I$, on a

\[
 \uline{\alpha_{i_0}},\beta_{i_0},\uline{\widetilde u_{i_0}},u_{i_0}\in TC
\]

(i.e. $\alpha_{i_0}$ \unsure{\&} $\widetilde u_{i_0}\in TC$, ce qui implique $\beta_{i_0}\in TC$ puis \inserted{$u_{i_0}\in TC$} par \unsure{cochgt.|cochgt$^{\rm t}$.} de base et par comp.) et ces conditions impliquent \inserted{$4^\circ)$} que $u_I:X_I\to Y_I$ est dans $TC$.

Voyons comment \deleted{\unsure{ou dans|on dé}} quelle mesure \unsure{le|ce} corollaire se déduit de la proposition. On a bien sûr

\diagram{\begin{tikzcd}[ampersand replacement=\&,column sep=small,row sep=small]
 \&3^\circ\arrow[dl,Rightarrow]\arrow[dr,Rightarrow]\&\\
 1^\circ\&\&2^\circ
\end{tikzcd}}

tautologiquement. D’ailleurs, la prop. appliquée à $TC$ \unsure{au lieu de} $C$ montre qu’on a $1^\circ\Rightarrow 3^\circ\mathbin{\&}4^\circ$, car en l’appliquant aux ens $I_{<i_0}$ on trouve que $\alpha_{i_0}\in TC$ $\forall i_0\in I$. Donc la proposition 2 implique

\[
 (1^\circ\Longleftrightarrow3^\circ)\Longrightarrow 2^\circ(\mathbin{\&}4^\circ).
\]
% archive sheet number: 20

\scanpage{22}
\mspage{9}

Donc pour prouver $2^\circ\Rightarrow3^\circ$, il \unsure{revient|reviendra} à prouver que $2^\circ$ \deleted{implique $4^\circ$} suffit à impliquer $4^\circ$ (et ceci quel que soit $I$ \inserted{\unsure{ordonné|considéré}} loc. fini) --- car en appliquant ceci aux $I_{<i_0}$, on trouve que $2^\circ$ implique $\alpha_{i_0}\in TC=C\cap W$ (non seulement $\alpha_{i_0}\in C$, selon la prop 2), donc $\beta_{i_0}\in TC\subset W$, donc $\widetilde u_{i_0}\in W$ par saturation, donc $\widetilde u_{i_0}\in TC$, d’où $3^\circ$.

Notons d’autre part que le b) de la prop.~2 \deleted{\unsure{n’est autre}} est conséquence \unsure{du|de} a), appliqué \deleted{\unsure{aux}} à la restriction de \unsure{d$^{\circ}$.|$u_\bullet$} à $I_{<i_0}$. Au total, il y a donc à prouver \unsure{seulement|seulem$^{\rm t}$} les deux choses

\[
\left\{\begin{aligned}
 \text{a)}\;&u_{\deleted{\unsure{I}}\bullet}\in C_I\Longrightarrow u_I\in C\\
 \text{b)}\;&u_\bullet\in C_I\cap W\Longrightarrow u_I\in C\cap W\quad(\overset{\text{déf}}{=}TC)\\
 &\hspace{5em}\inserted{[\text{et }u_\bullet\in(TC)_{\unsure{W|w}}]}
\end{aligned}\right. \tag*{(*)}
\]

Nous le prouverons par réduction \deleted{\unsure{à la}} aux cas particuliers de la prop.~1, déjà prouvée. J’ai d’ailleurs \unsure{aussi|ici} envie de prouver le

\uline{Corollaire 2} Les hypothèses \inserted{$\unsure{\mathcal M|\mathscr M},C$ (\deleted{\unsure{et}} \unsure{évt.|évent.} $W$),} sur $X_\bullet,Y_\bullet\in\mathcal M(I^\circ)_C$ et sur $u_\bullet\in C_I$, \inserted{\unsure{enfin|et} sur $I$ loc. fini,} sont celles de la prop.~2. \deleted{Soient} \inserted{Pour \unsure{tte} \inserted{partie} $J$ de $I$, soit} \deleted{$I'\subset I''$ deux \unsure{ouvertes|parties}}

\[
 \deleted{\unsure{X_J}}u_J\overset{\text{déf}}{=}\underset{i\in J}{\underrightarrow{\lim}}\,u_i:
 X_J\overset{\text{déf}}{=}\underset{\unsure{i'|i}\in J}{\underrightarrow{\lim}}\,X_i
 \mathbin{\unsure{-|\to}} Y_J\overset{\text{déf}}{=}\underset{i\in J}{\underrightarrow{\lim}}\,Y_i .
\]

Soient $I'\unsure{\subseteq|\subset} I''$ deux parties \uline{ouvertes} de $I$, et considérons \unsure{le|les} \unsure{diagramme|diagrammes}
% archive sheet number: 21

\scanpage{23}
\mspage{10}

\diagram{\begin{tikzcd}[ampersand replacement=\&,column sep=large]
 \&X_{I'}\arrow[r,"\tau_X"]\arrow[d,"\alpha=u_{I'}"']\arrow[dr,phantom,"\text{cocart}" description]
 \&X_{I''}\arrow[d,"\beta"']\arrow[dd,bend left=43,"u_{I''}"]\\
 (*)\&Y_{I'}\arrow[r,"\widetilde\tau_Y"]\arrow[dr,dashed,no head,"\tau_Y"']
 \&\widetilde Y_{I''}\arrow[d,"\widetilde u"'] \&\text{(défini comme \unsure{somme} \unsure{amalgamée})}\\
 \&\&Y_{I''}
\end{tikzcd}}

Alors on a \deleted{$\unsure{\alpha},\beta,\widetilde u$.}

\[
\left\{\begin{aligned}
 &u_{I'},u_{I''}\in C\quad (\text{i.e. pour tt }J\text{ ouvert \unsure{dans|de} }I,\!\\
 &\qquad \inserted{[\text{donc aussi }\beta\in C]}\quad u_J:X_J\mathbin{\unsure{-|\to}}Y_J\text{ est dans }C)^{*}\\
 &\text{et }\deleted{\beta,\unsure{\widetilde u}}\ \tau_X,\tau_Y\in C\quad(\text{i.e. pour }I'\subset I''\\
 &\qquad\text{inclusion d’ouverts, et }X\ C\text{-cofibrant dans}\\
 &\qquad\unsure{\mathcal M|\mathscr M}(I^\circ),\ \unsure{\ldots|\text{on a}}\ X_{I'}\mathbin{-}X_{I''}\text{ dans }C)^{**}\\
 &(\text{\unsure{donc aussi|d’où encore} }\deleted{\unsure{\beta},}\widetilde\tau_Y\in C)\\
 &\deleted{\text{\unsure{et de plus}}}\ \text{et \unsure{enfin}, même}\\
 &\boxed{\widetilde u\in C}\,.
\end{aligned}\right.
\]

% slanted
\margin{$^{*}$c’est bien sûr contenu dans prop.~2, appliqué à $u_\bullet|J$}

% slanted
\margin{C’est pour cette assertion $\tau_X,\deleted{\unsure{\tau_Y}}\in C$ (d’où $\widetilde\tau_Y,\tau_Y\in C$ moyennant $\widetilde u\in C$) qu’on a besoin de l’hypothèse $X_\bullet\in\mathcal M(I^\circ)_C$. Mais $\widetilde u\in C$ vrai \unsure{sans} cette hypothèse}

\deleted{\unsure{L}} \uline{NB} Le cas de la prop.~2 \inserted{sous la forme cas a) (page 9),} \deleted{est \unsure{obtenu}}, (i.e. \unsure{c’est} le fait que \inserted{les} $(u_{I'},u_{I''}\ldots$ sont dans $C$) est obtenu en prenant $I''=I,\ I'=\varnothing$, le diagramme devient

\diagram{\begin{tikzcd}[ampersand replacement=\&,row sep=small]
 \varnothing\arrow[r]\arrow[d,no head,"\mathrm{id}"']\arrow[dr,phantom,"\text{cocart}" description]\& X_I\arrow[d,"\mathrm{id}"]\\
 \unsure{\varnothing}\arrow[r]\&X_I\arrow[d,"\widetilde u="' ,"u_I"]\\
 \&Y_I
\end{tikzcd}}

et la conclusion est donc $u_I\in C$. \inserted{$^{*}$} La \uline{définition} de $C_I$, i.e. la condition $u\in C_I$, est \unsure{retrouvée} en exigeant $\widetilde u\in C$ \inserted{\unsure{dans les cas|du cas}} ci-dessus, dans les cas $I''=I_{/i_0}$, $I'=I_{<i_0}$. Il est remarquable

% margin note written sideways
\margin{$^{*}$ Le cas $(\deleted{\unsure{X}}Y_{I'}\xrightarrow{\tau_{\unsure{Y|X}}}\deleted{\unsure{X}}Y_{I''})\in C$ est obtenu en prenant $X_\bullet=\varnothing_\bullet$, on aura
\diagram{\begin{tikzcd}[ampersand replacement=\&,column sep=tiny,row sep=tiny]
 \varnothing\arrow[r,"\approx","\tau_X"']\arrow[d]\arrow[dr,phantom,"\text{cocart.}" description]\&\varnothing\arrow[d,no head]\\
 Y_{I'}\arrow[r,"\mathrm{id}"]\arrow[dr,"\tau_Y"']\&\widetilde Y_{I''}=Y_{I'}\arrow[d,"\widetilde u={}"' ,"\tau_Y"]\\
 \&Y_{I''}
\end{tikzcd}}}
% archive sheet number: 22

\scanpage{24}
\mspage{11}

que cette définition de $C_I$ donne une conclusion aussi forte que $\widetilde u\in C$ pour \uline{tt} couple $I'\subset I''$ de sous-ens. ouverts de $I$.

C’est finalement sous cette forme qu’on va prouver la prop.~2 --- c’est \uline{\unsure{ça|là}} à présent qui me semble \uline{le} bon énoncé dans cette direction. Avec le complément, bien sûr, que si de plus $u_\bullet\in W_{I^\circ}$ i.e. les $u_i$ sont $\in W$, alors $u_\bullet$ est $\unsure{=}{}\in(TC)_I$, \unsure{i.e.|voire} aussi bien que \deleted{\unsure{les}} les \uline{$u_{I'}$ sont $\in TC$} (d’où il \unsure{résultera|résulte} \deleted{\unsure{que}} donc que $\widetilde u$ du cas p.~10 est aussi dans $TC$).

On va procéder par construction transfinie. On peut supposer bien sûr que $I''=I$. On va \inserted{peut $=$} d’abord prouver le cas particulier \deleted{\unsure{où}} \deleted{il existe} où $I'=\varnothing$, i.e. prouver \deleted{que} \unsure{montrer|aussi} que $u_I\in C$. On \deleted{va proc\unsure{.}} construit par réc. transfinie une suite \inserted{str.} croissante d’ouverts

\[
 I_\alpha\subset I,
\]

aboutissant à $I_{\alpha_0}=I$,

en posant

\[
\left\{\begin{aligned}
 I_0&=\varnothing\\
 I_\alpha&=\bigcup_{\beta<\unsure{\alpha|2}}I_{\unsure{\alpha|\beta}}
       &&\text{si }\alpha\text{ ordinal limite}\\
 I_{\alpha\unsure{+|{}}1}&=I_\alpha\cup I_{/s_\alpha}
       &&\text{si }I_\alpha\ne I,\text{ et en choisissant}\\
 &&&s_\alpha\in I\smallsetminus I_\alpha.
\end{aligned}\right.
\]
% archive sheet number: 23

\scanpage{25}
\mspage{12}

(On finit par avoir $I_{\alpha_0}=I$, bien sûr.) Les \deleted{\unsure{pose}}

\[
\begin{aligned}
 X_\alpha&=X_{I_\alpha}=\underset{i\in I_\alpha}{\underrightarrow{\lim}}\,X_i
       &&\unsure{u}\\
 Y_\alpha&=Y_{I_\alpha}=\underset{i\in I_\alpha}{\underrightarrow{\lim}}\,Y_i\\
 u_\alpha&=u_{I_\alpha}:X_\alpha\mathbin{\unsure{-|\to}}Y_\alpha,
\end{aligned}
\]

\unsure{forment|hom} \unsure{des|de} systèmes inductifs transfinis. On aura

\[
 \deleted{\unsure{u}}X_{\alpha_0}=X_I,\ Y_{\alpha_0}=Y_I,\ u_{\alpha_0}=u_I,
\]

et on doit \inserted{donc} prouver que $u_{\alpha_0}\in C$. Pour ceci on applique la prop.~1, \deleted{\unsure{elle}} \deleted{\unsure{il nous|on aura}} \inserted{suffit} donc de prouver que $(u_\alpha)\in C_{I_{\alpha_0+1}}$,

\[
 [\,\text{où }I_{\alpha_0+1}=\{\text{ordinaux }i\mid i<\alpha_0+1\ \text{i.e. }i\leq\alpha_0\}\,]
\]

\deleted{est dans $C$} Donc il faut prouver que les flèches $\widetilde u_{i_0}$ du cas p.~1 (où cette fois les $i,i_0$ sont des ordinaux, $i_0\leq\alpha_0$) sont dans $C$. Je vais distinguer les deux cas où $\alpha=i_0$ est ordinal \deleted{\illegible{1}} limite, \deleted{\unsure{qui est|quelque} \unsure{trivial}} \inserted{trivial car alors} $X_{i_0}\xleftarrow{\sim}\underset{i<i_0}{\lim}X_i$, $Y_{i_0}\xleftarrow{\sim}\underset{i<i_0}{\underrightarrow{\lim}}Y_i$, et $\widetilde u_{i_0}=\mathrm{id}$, et le cas $\alpha$ non \deleted{\unsure{limite}}, \unsure{i.e.|ici}.

\diagram{\begin{tikzcd}[ampersand replacement=\&,column sep=huge]
 X_{I_{\alpha-1}}\arrow[r]\arrow[d,"u_{\alpha-1}"']\&X_{I_\alpha}\arrow[d,"\beta_\alpha"]\arrow[dd,bend left=46,"u_\alpha"]\\
 Y_{I_{\alpha-1}}\arrow[r,no head]\&\widetilde Y_{I_\alpha}\arrow[d,"\widetilde u_\alpha"]\\
 \&Y_{I_\alpha}
\end{tikzcd}}
% archive sheet number: 24

\scanpage{26}
\mspage{13}

à prouver que $\widetilde u_\alpha\in C$. Lorsque on suppose de plus $u_\bullet\in W_{I^\circ}$ i.e les $u_i\in W$, il faut de plus prouver que $\widetilde u_\alpha\in TC=C\cap W$. La prop.~2 et son corollaire 1 résulteront de là. Pour prouver $\widetilde u_\alpha\in C$ \deleted{\unsure{i.e.}} resp. $\widetilde u_\alpha\in TC$, on \deleted{utilise sa \unsure{proc.}} peut supposer $I_\alpha=I$ (quitte à changer de notation). On pose donc

\[
\begin{aligned}
 I_\alpha&=I,\ I_{\alpha-1}=I'\\
 I_{/s_{\alpha-1}}&=I_0,\ I_0\cap I'=I'_0,
\end{aligned}
\]

de sorte qu’on a un diagramme commutatif d’inclusions ouvertes

\diagram{\begin{tikzcd}[ampersand replacement=\&]
 I_0\cap I'=I'_0\arrow[r,hook]\arrow[d,hook]\&I_0\arrow[d,hook]\\
 I'\arrow[r,hook]\& I=I_0\cup I'\, .
\end{tikzcd}}

On utilisera seulement le fait que $I_0$, donc $I'_0$, est fini, i.e que \inserted{dans la \unsure{config.|cond.} ($I'\subset I$),} $I$ peut s’écrire sous la forme $I=I'\cup I_0$, avec $I_0$ un ouvert \uline{fini}, pour en conclure que dans le diagramme

\diagram{\begin{tikzcd}[ampersand replacement=\&,column sep=large]
 X_{I'}\arrow[r,"{\tau^X_{I,I'}}"]\arrow[d,"u_{I'}"']\arrow[dr,phantom,"\text{cocart}" description]
 \&\deleted{\unsure{I}}X_{\deleted{\unsure{1|{}}}I}\arrow[d,"{\unsure{v|u}_{I',I}}"]\arrow[dd,bend left=42,"u_I"]\\
 Y_{I'}\arrow[r]\arrow[dr,"{\tau^{Y}_{I,I'}}"']\&\widetilde Y_I\arrow[d,"{\widetilde u=u_{I',I}}"]\\
 \&Y_I
\end{tikzcd}}

on a $\widetilde u_{\deleted{\unsure{I}}}\in C$ resp $\widetilde u\in TC=C\cap W$.
% archive sheet number: 25

\scanpage{27}
\mspage{14}

Il faut des notations générales, \deleted{pour} d’une part pour une inclusion $J\subset K$ d’ouverts de $I$, donnant lieu au diagramme

\diagram{\begin{tikzcd}[ampersand replacement=\&,column sep=large]
 \& X_J\arrow[r,"{\tau^{X}_{J,K}}"]\arrow[d,"u_J"']\arrow[dr,phantom,"\text{cocart}" description]
 \& X_K\arrow[d,"{\widetilde u_{J,K}}"']\arrow[dd,bend left=50,dash pattern=on 9mm off 1.2mm on .7mm off 1.2mm on .7mm off 1.2mm on .7mm off 1.2mm on .7mm off 1.2mm on .7mm off 1.2mm on .7mm off 1.2mm on .7mm off 1.2mm on .7mm off 1.2mm,"u_K"]\\
 (1)\&Y_J\arrow[r]\arrow[dr,dashed,no head,"{\tau^{Y}_{J,K}}"']\&Y_{J,K}\arrow[d,"{u_{J,K}}"']\&(\text{\unsure{anciennt|anciennement} }\widetilde Y_{\unsure{K|k}})\\
 \&\&Y_K\&{,}
\end{tikzcd}}

et pour \deleted{deux} inclusion de trois ouverts emboîtés $J\subset K\subset L$ de $I$,

\diagram{\begin{tikzcd}[ampersand replacement=\&,column sep=large]
 \&X_J\arrow[r]\arrow[d,"u_J"']\arrow[dr,phantom,"\text{cocart}" description]
 \&X_K\arrow[r]\arrow[d,"{\widetilde u_{J,K}}"']\arrow[dr,phantom,"\text{cocart}" description]
 \&X_L\arrow[d,"{\widetilde u_{J,L}}"']\arrow[dd,bend left=33,"{\widetilde u_{K,L}}"]\\
 (2)\&Y_J\arrow[r,no head]\&Y_{J,K}\arrow[r]\arrow[d,"{u_{J,K}}"']\arrow[dr,phantom,"\text{cocart}" description]
 \&Y_{J,L}\arrow[d,"{u^L_{J,K}}"']\arrow[dd,bend left=48,"{u_{J,L}}"]\\
 \&\&\deleted{\unsure{Y_{J}}}Y_K\arrow[r]\& Y_{K,L}\arrow[d,"{u_{K,L}}"']\\
 \&\&\&Y_L\ ,
\end{tikzcd}}

donnant lieu à une nouvelle flèche $u^{L}_{\unsure{J|I},K}$, \unsure{déduisant|déduite} de $u_{J,K}$ par le \unsure{coch.|ch.} de base
% archive sheet number: 26

\scanpage{28}
\mspage{15}

$Y_{J,K}\to\deleted{\unsure{X}}Y_{J,L}$. Nous utiliserons ce diagramme dans les deux cas où $(J\subset K\subset L)$ est l’un de deux triples

\[
 \left\{\begin{aligned}
 I'_0&\subset I_0\subset I\\
 I'_0&\subset I'\subset I
 \end{aligned}\right.
\]

\inserted{issus} du carré cocartésien (\deleted{\unsure{de la}} cat) (\deleted{\unsure{ou}} carré MV)

\diagram{\begin{tikzcd}[ampersand replacement=\&]
 \&I'_0\arrow[r,hook,no head]\arrow[d,hook]\&I_0\arrow[d,hook]\\
 (3)\&I'\arrow[r,hook]\&I\, .
\end{tikzcd}}

Il faut insérer les deux diagrammes correspondants \unsure{du|au} type (2) dans un même diagramme \unsure{!|?} On notera à ce propos que les carrés \inserted{\unsure{de|dans} $\unsure{\mathcal M|\mathscr M}$,} déduits de (3) en lui appliquant \deleted{$\unsure{J}$} $J\mathbin{\unsure{\mapstochar\relbar|\mapsto}}X_J$ et $J\mathbin{\mapstochar\relbar}Y_J$

\diagram{$\left.\vcenter{\hbox{\begin{tikzcd}[ampersand replacement=\&,column sep=small,row sep=small]
 \&X_{I'_0}\arrow[rr]\arrow[dr]\arrow[dd,dashed,"u_{I'_0}"]\&\&X_{I_0}\arrow[dr]\arrow[dd,dashed,no head,"u_{I_0}"]\&\\
 (4_x)\&\&X_{I'}\arrow[rr]\arrow[dd,dashed,"u_{I'}"']\&\&X_I\arrow[dd,dashed,"u_I"']\\
 \&Y_{I'_0}\arrow[rr]\arrow[dr]\&\&Y_{I_0}\arrow[dr]\&\\
 (4_y)\&\&Y_{I'}\arrow[rr,no head]\&\&Y_I
\end{tikzcd}}}\right\}\ (4)$}

sont cocartésiens dans $\mathcal M$. (Transitivité des $\underrightarrow{\lim}\ldots$).
% archive sheet number: 27

\scanpage{29}
\mspage{16}

On a un hom. \inserted{$(\deleted{u_{I'_0},u_{I_0}},u_{I'},u_I)$} déduit de $u$ du carré $4_x$ dans le carré $4_y$, et cet homomorphisme lui-même s’insère dans un diagramme, incluant les deux autres déjà signalés de (2). Voici le diagramme pertinent :

\diagram{\begin{tikzcd}[ampersand replacement=\&,column sep=small,row sep=small]
 \&\&X_{I_0}\arrow[rrr,Rightarrow]\arrow[dd,dashed,no head]\&\&\&X_I\arrow[dd]\&\\
 \&X_{I'_0}\arrow[ur]\arrow[rr]\arrow[dd,no head]\&\&X_{I'}\arrow[urr]\arrow[dd,no head]\&\&\&\\
 \&\&Y_{I'_0,I_0}\arrow[rrr,dash pattern=on 2mm off 3mm on 8mm off .2mm,no head]\arrow[dd,very thick,"{\unsure{u^{\wedge}_{I'_0,I_0}|u_{I'_0,I_0}}}"]\&\&\&
 Y_{I'_0,I}\arrow[ddr,very thick,"{u^I_{I'_0,I_0}}"]\arrow[ddd,thick]\&\\
 (5)\&Y_{I'_0}\arrow[ur]\arrow[rr]\arrow[dr,dash dot]\arrow[dddrr,dashed]\&\&
 Y_{I'_0,I'}\arrow[urr,thick]\arrow[ddd]\&\&\&\\
 \&\&Y_{I_0}\arrow[rrrr,thick]\arrow[ddrrrr,dashed]\&\&\&\&Y_{I_0,I}\arrow[dd,thick,"{u_{I_0,I}}"]\\
 \&\&\&\&\&Y_{I',I}\arrow[dr,very thick,"{u_{I',I}}"']\&\\
 \&\&\&Y_{I'}\arrow[urr]\arrow[rrr,dashed]\&\&\&Y_I
\end{tikzcd}}

% slanted
\margin{diagramme \unsure{imbitable|imbittable|indébrouillable|indéchiffrable}, à remplacer par diagramme (6) ci-dessous}

\deleted{\unsure{Dans}} Le cube supérieur est le cube \inserted{de $\unsure{\mathcal M|\mathscr M}$,} à 6 faces cocartésiennes déduit des trois flèches de $X_{I'_0}$ dans $X_{I_0},X_{I'},Y_{I'_0}$. \inserted{On aura besoin de savoir que le \uline{bas} du cube est un carré cocartésien.} La partie antérieure du diagramme est le \unsure{dessin} de \unsure{(2)|(1)} correspondant aux inclusions $I'_0\subset I'\subset I$, la partie derrière \deleted{est} l’analogue pour \unsure{l’inclusion|les inclusions} $I'_0\subset I_0\subset I$. Les deux autres carrés verticaux, en dessous du cube supérieur, sont également cocartésiens, ainsi que le carré oblique en pointillés, lequel n’est autre que le bas $4_y$ du diagramme (4). De ces \unsure{caractères|conditions} du diagramme 5 résulte \deleted{\unsure{que}} aisément (\unsure{quoique|puisque}
% archive sheet number: 28

\scanpage{30}
\mspage{17}

j’ai bien mis une bonne journée de \unsure{gribouillis} en \unsure{barbouillant|tournant|tournaillant} là-dedans, avant de m’en rendre compte enfin\ldots) que le carré en losange à traits pleins

\diagram{\begin{tikzcd}[ampersand replacement=\&,column sep=small]
 \&Y_{I'_0,I}\arrow[dl]\arrow[dr]\&\\
 Y_{I',I}\arrow[dr]\&\&Y_{I_0,I}\arrow[dl,no head]\\
 \&\deleted{\unsure{Y|X}}Y_I\&
\end{tikzcd}}

est également cocartésien. On le voit \unsure{trivialem$^{\rm t}$|trivialement} sur le diagramme déployé

\diagram{\begin{tikzcd}[ampersand replacement=\&,row sep=small,column sep=tiny]
 \&\&\&Y_{I'_0}\arrow[dl]\arrow[dr,thick]\arrow[dd,phantom,"\text{cocart}" description]\&\&\\
 \&\&Y_{I'_0,I'}\arrow[dl,"{\deleted{\unsure{u_{I'_0,I'}}}}"']\arrow[dr]\arrow[dd,phantom,"\text{cocart}" description]\&\&
 Y_{I'_0,I_0}\arrow[dl]\arrow[dr,"{u_{I'_0,I_0}}"]\arrow[dd,phantom,"\text{cocart}" description]\&\\
 (6)\&Y_{I'}\arrow[dr]\&\&Y_{I'_0,I}\arrow[dl,thick,"{\deleted{\unsure{u^I_{I'_0,I_0}|u^I_{I'_0,I'}}}}"]\arrow[dr,"{u^I_{I'_0,I_0}}"]\arrow[dd,phantom,"\text{cocart}" description]\&\&Y_{I_0}\arrow[dl]\\
 \&\&Y_{I',I}\arrow[dr,very thick,"{\unsure{u|v|\widehat u}_{I',I}}"]\&\&Y_{I_0,I}\arrow[dl,"{\deleted{\unsure{u_{I_0,I}}}}"]\&\\
 \&\&\&{?}\simeq Y_I=Y_{I'}\amalg_{Y_{I'_0}}Y_{I_0}\ ,\&\&
\end{tikzcd}}

qui est finalement le diagramme pertinent (le diagramme (5) ne servant à rien, c’est un mauvais bout par lequel j’ai pris les choses). On voit sur le diagramme \inserted{(6)} que pour prouver $u_{I',I}\in C$ resp. $TC$, il
% archive sheet number: 29

\scanpage{31}
\mspage{18}

suffit de prouver $u_{I'_0,I_0}\in C$ resp. $TC$. Mais c’est là une assertion relative à l’inclusion \deleted{de \unsure{catégories}} \inserted{d’ens. ordonnés} $I'_0\subset I_0$ \uline{finis}. On est donc réduit à prouver le lemme suivant, où $I_0$ devient $I$ :

\uline{Lemme} : Soit $I$ un ens. ordonné fini, \deleted{$\unsure{I'}$ un ouvert de $I$,} $u_\bullet:X_\bullet\mathbin{\unsure{-|\to}}Y$ une flèche dans $\unsure{\mathcal M|\mathscr M}(I^\circ)$, où $\mathcal M$ est une catégorie munie d’un ens. $C\subset\mathrm{Fl}(\mathcal M)$ de flèches, tel que $C$ soit stable par coch. de base et par \inserted{\unsure{(*)}} composition, et contienne les iso.

On suppose $u_\bullet\in C_I$ i.e. que pour $i\in I$, la flèche $u_{I_{<i},I_{/i}}$ est dans $C$. Alors pour toute partie ouverte $I'$ de $I$, la flèche $u_{I',I}$ est dans $C$.

% slanted
\margin{\uline{NB} Bien sûr, $u_{J,K}\in C$ si $J\subset K$ sont deux ouverts comme on voit en appliquant le lemme à $K$ \unsure{ald} $I$, $\unsure{\to|\checkmark}$}

\inserted{\uline{\unsure{Cor.}}} Si de plus $W$ \inserted{$\subset\mathrm{Fl}\,\mathcal M$} est un localiseur modérément \unsure{ramifié|saturé}, tel que \deleted{\unsure{alors}} $TC\overset{\text{déf}}{=}W\cap C$ soit stable par coch. de base, \deleted{\unsure{alors est}} (il l’est de ttes façons par composition, et contient les iso), alors on a

\[
 C_I\cap W_{I^\circ}=(TC)_I,
\]

donc si $u_\bullet\in C_I$ (cf ci-dessus) et les $u_i\in W$, alors $u_\bullet\in(TC)_{\unsure{I}}$, donc \deleted{\unsure{par}} appliquant ce qui est dit pour $C$ à $TC$, on a $u_{I',I}\in TC$ pour \unsure{ttes|tte}
% archive sheet number: 30

\scanpage{32}
\mspage{19}

partie ouverte $I'$ de $I$. \deleted{\unsure{Si}}

On prouvera en $=$ \deleted{\unsure{temps}} le lemme et son corollaire, en procédant par récurrence sur $\mathrm{card}.\ I=n$. Si $n\unsure{\leq|=}0$ i.e. $I=\varnothing$, les énoncés sont triviaux. Par transitivité, en utilisant les diagrammes du type \unsure{(1?)|(1)?|(2?)}, on est ramené au cas où $I=I'\cup\{i_0\}$, où $i_0\in I\smallsetminus I'$. On pose \inserted{encore} $I_0=I_{/\unsure{i|l}_0}$, $I'_0=I_0\cap I'$, et le raisonnement qu’on vient de faire nous montre qu’il suffit de prouver que $u_{I'_0,I_0}\in C$ resp. $\in TC$. Si $I'_0\ne I_{\deleted{\unsure{0|g}}}$, cela résulte de l’hyp. de récurrence. Sinon, $i_0$ est un objet final de $I$, et $I'=I_{<i_0}$, et alors $u_{I',I}\in C$ \deleted{résulte} est inclus dans la définition de $C_I$. \deleted{\unsure{Soit}} Reste à prouver \inserted{donc} que $u_{I',I}\in TC$ si les $u_i\in W$. On a le diagramme

\diagram{\begin{tikzcd}[ampersand replacement=\&,column sep=large]
 \unsure{X}_{I'}\arrow[r]\arrow[d,"u_{I'}"']\&X_I\arrow[d,"{\widetilde u_{I',I}}"]\arrow[dd,bend left=45,dashed,"u_I"]\\
 Y_{I'}\arrow[r]\&Y_{I',I}\arrow[d,"{u_{I',I}}"]\\
 \&Y_I\ ,
\end{tikzcd}}
% archive sheet number: 31

\scanpage{33}
\mspage{20}

et l’hyp. de récurrence implique $u_{I'}\in TC$, donc $\widetilde u_{I',I}\in TC$ par coch. de base. D’autre part la flèche $u_I:X_I\mathbin{\unsure{-|\to}}Y_I$ s’identifie à $X_{i_0}\mathbin{-}Y_{i_0}$ ($i_0$ étant élément final de $I$), elle est donc dans $W$. Donc par saturation, $u_{I',I}\in W$, donc $u_{I',I}\in TC$, \unsure{qed.|cqfd.}

Ainsi, la prop.~2 est prouvée, ainsi que son cor.~1 --- et on voit que si dans la prop 2 et son cor., on suppose $I$ \uline{fini}, alors il est inutile dans les hyp pour $C$, resp. $TC=C\cap W$, de supposer la stabilité par composition \uline{transfinie}, la composition ordinaire suffit. Mais il reste \unsure{à} prouver le corollaire 2, que je viens de prouver dans le cas \deleted{\unsure{de}} $I$ \uline{fini}, sans avoir à faire appel à une hypothèse ``transfinie''. Il faut maintenant prouver que c’est \unsure{aussi|encore} OK \deleted{\unsure{si}} dans le cas loc fini, en supposant $C$ stable par \unsure{composition|compositions} \unsure{transfinie|transfinies} \unsure{!} (Il n’y a plus à revenir sur le cas de $W$, \deleted{\unsure{on}} on sait déjà \inserted{maintenant} que $C_I\cap W_{I^\circ}=(TC)_I$, par le cor.~1 déjà prouvé.)
% archive sheet number: 32

\scanpage{34}
\mspage{21}

La démonstration est \inserted{essentiellement} la $=$ que pour le cas particulier déjà traité $I'=\varnothing,\ I''=I$, et on aurait pu aussi bien rédiger la démonstration dans ce cadre, à l’aide d’une suite \inserted{str. croissante,} transfinie d’ouverts \inserted{de $I$} $I_\alpha$, avec $\boxed{I_0=I'}$, $I_{\alpha_0}=I''(=I)$.

\deleted{On peut d’ailleurs ramener le cas général au cas déjà traité, en introduisant, quand $I'\subset I(=I'')$ \unsure{fixés|fixé}, \deleted{\unsure{des}} une nouvelle flèche
\[
 u_\bullet:X'_\bullet\longrightarrow Y'_\bullet,
\]
en posant
\[
 X'_i=\underset{j\in I'\cup I_{/i}}{\underrightarrow{\lim}}\,X_j,\quad
 Y'_i=\underset{j\in I'\cup I_{/i}}{\underrightarrow{\lim}}\,Y_j
\]
et \deleted{en travaillant dans} $\unsure{\mathcal M|\mathscr M}'={}_{\unsure{x|X}_{I'}}$
\[
 X'_{\varnothing}=X_{I'},\quad Y'_{\varnothing}=Y^{\unsure{4|\leftarrow}}_{I'}
\]
$\unsure{X'}$}

et

\[
\begin{aligned}
 I_\alpha&=I_{\alpha-1}\cup I_{/s_{\alpha-1}}
 &&\text{si }\alpha\text{ non limite, }\inserted{I_{\alpha-1}\ne I,}\ \deleted{\text{pour}}\ \text{choisissant}\\
 &&&s_{\alpha-1}\in I\smallsetminus I_{\alpha-1}\\
 I_\alpha&=\bigcup_{\beta<\alpha} I_\beta\quad .
 &&\text{si }\alpha\text{ limite.}
\end{aligned}
\]

et posant encore

\[
 X_\alpha=X_{I_\alpha},\quad Y_\alpha=Y_{I_\alpha},\quad u_\alpha=u_{I_\alpha}:X_\alpha\mathbin{\unsure{-|\to}}Y_\alpha.
\]

Soit $\alpha_0$ tel que \deleted{$\unsure{\alpha}$} $I_{\alpha_0}=I$ (indice terminal de la suite) ; \deleted{\unsure{alors} \unsure{il} \unsure{prouv.}} il faut prouver que \unsure{dans|les|la}

\diagram{\begin{tikzcd}[ampersand replacement=\&]
 X_0\arrow[r]\arrow[d,no head]\&X_{\alpha_0}\arrow[d,no head]\\
 Y_0\arrow[r]\&Y_{0,\alpha_0}\arrow[d,"{u_{0,\alpha_0}}"]\\
 \&Y_{\alpha_0}
\end{tikzcd}}
% archive sheet number: 33

\scanpage{35}
\mspage{22}

on a

\[
 u_{0,\alpha_0}\in C.
\]

Compte tenu du complément marginal à proposition 1 (qui se prouve de la $=$ façon, \unsure{\ldots}), on est ramené à prouver que \deleted{\unsure{des|les}} \inserted{l’hom. des} systèmes ordinaux envisagés est dans $C_{\unsure{\mathbf I|\mathbb I}_{\alpha_0+1}}$ \deleted{\unsure{a}} (où $\mathbf I_{\alpha_0+1}$ est l’ens des ordinaux $\leq\alpha_0$). Or \unsure{qu’il} en \unsure{soit ainsi} se \unsure{prouve} \inserted{\unsure{exact$^{\rm t}$|exact}} comme tantôt, à cela près \deleted{\unsure{qu’on}} \deleted{\unsure{que} \unsure{la} \unsure{vérific.}} \inserted{qu’on doit \inserted{\unsure{à présent}} faire} une vérification \unsure{du} cas p.~1 ($\widetilde u_{\unsure{v|\nu}_0}\in C$) pour $v_0=0$. Alors le diagramme \unsure{cas|de} p.~1 est

\diagram{\begin{tikzcd}[ampersand replacement=\&]
 \varnothing_{\deleted{\unsure{0|I}}}\arrow[r]\arrow[d,no head,"{\mathrm{iso}}"]\&X_0=X_{I'}\arrow[d,"{\mathrm{iso}}"',"{\mathrm{id}}"]\\
 \varnothing_{\deleted{\unsure{0|I}}}\arrow[r]\&X_0=X_{I'}\arrow[d,"{u_0=u_{I'}}"]\\
 \&Y_0=Y_{I'}
\end{tikzcd}}

et il faut donc vérifier $u_0=u_{I'}\in C$. Or ceci est contenu dans la prop.~2 elle-$=$, déjà prouvée, qu’on doit donc utiliser. Le reste se prouve comme avant. Ainsi, le cor.~2 de prop.~\unsure{2} est également prouvé

\uline{Commentaires} : Il y a un malaise, car les prop.~1 et prop.~2 \deleted{\unsure{sont}} (ainsi que leurs corollaires) sont \unsure{essentielt|essentiellem$^{\rm t}$} les $=$ \unsure{conclusions|conclusion}, avec des hypothèses \inserted{sur $I$} qui sont
% archive sheet number: 34
